
\section{$\OO$ 模的 $\mathrm{Ext}$ 群与 $\mathrm{Ext}$ 层}

回忆一下：对任意\trRingedSpace{} $X$，范畴 $\mathit{Mod}_{\OO_X}$ 有足够多内射对象（\mbox{定理~23.4.1}）。于是，对 $X$ 上任意 $\OO_X$-\trModule{} $\mathscr{F}$，我们可以定义
\[
\mathrm{Ext}^i_X(\mathscr{F},\cdot)\colon\mathit{Mod}_{\OO_X}\longrightarrow\mathit{Mod}_{\Gamma(X,\OO_X)}
\]
为 $\Hom_X(\mathscr{F},\cdot)$ 的第 $i$ 个右导出函子，并且 $\mathrm{Ext}^i_X(\mathscr{F},\cdot)$ 有相应的\trLongExactSequence。我们类似地定义层版本
\[
\mathscr{E}xt^i_X(\mathscr{F},\cdot)\colon\mathit{Mod}_{\OO_X}\longrightarrow\mathit{Mod}_{\OO_X}
\]
为 $\mathscr{H}om_X(\mathscr{F},\cdot)$ 的第 $i$ 个右导出函子。在这两种情形下，当上下文清楚时，下标 $X$ 常被省略，尽管当讨论中涉及多于一个空间时，这样做可能很危险。（我们在 \S23.2.4 中见过 $A$-\trModule{}的 $\mathrm{Ext}$ 函子。）

\emph{警示：}目前尚不清楚 $\mathit{Mod}_{\OO_X}$ 是否有足够多\trProjectiveObj{}对象（事实上并不成立；见 \S23.4.8），所以我们不能把 $\mathrm{Ext}^i$ 定义成对左变元取导出函子。尽管如此，我们将会看到，它的表现就好像它\emph{是}一个导出函子——它「可用零调对象计算」，并且具有\trLongExactSequence（\mbox{注记~29.2.1}）。

\emph{又一警示：}在这个定义下，如果 $\mathscr{F}$ 与 $\mathscr{G}$ 是一个\trScheme{}上的\trQCSheaf，尚不清楚诸 $\mathscr{E}xt^i(\mathscr{F},\mathscr{G})$ 是否\trQC；而且，\mbox{练习~14.3.A(a)} 中的旁注指出，即便对 $i=0$，这也未必成立。但 \mbox{练习~29.2.F} 会让你安心：只要 $\mathscr{F}$ \trCoh{}，而 $X$ \trLocally\trNoetherian，一切安好。

\mbox{练习~23.5.A}（内射 $\OO_X$-\trModule{}限制到\trOpenSubset{} $\mathfrak{U}\subset X$ 上时，在 $\mathfrak{U}$ 上仍内射）有一系列有用的推论。

\textbf{29.2.A.（练习）}\quad 设 $\mathscr{I}$ 是一个内射 $\OO_X$-\trModule。证明：$\mathscr{H}om_{\OO_X}(\cdot,\mathscr{I})$ 是一个\trExact{}\trContravariantFunctor。（一个相关的事实：$\Hom_{\OO_X}(\cdot,\mathscr{I})$ 是\trExact{}的，由内射性的定义，\mbox{练习~23.2.C(a)}。）

\textbf{29.2.B.（练习）}\quad 设 $X$ 是一个\trRingedSpace，$\mathscr{F}$ 与 $\mathscr{G}$ 是 $\OO_X$-\trModule，而 $\mathfrak{U}$ 是一个\trOpenSubset。请描述一个\trCanonical\trIsomorphism
\[
\mathscr{E}xt^i_X(\mathscr{F},\mathscr{G})|_{\mathfrak{U}}\xrightarrow{\ \sim\ }\mathscr{E}xt^i_{\mathfrak{U}}\bigl(\mathscr{F}|_{\mathfrak{U}},\mathscr{G}|_{\mathfrak{U}}\bigr).
\]
提示：限制函子 $\mathit{Mod}_{\OO_X}\to\mathit{Mod}_{\OO_Y}$ 是\trExact{}的（参看\mbox{练习~2.7.D} 与\mbox{练习~2.7.E}）；限制保持 $\OO$-\trModule{}的内射性（\mbox{练习~23.5.A}）；以及 FHHF 定理（\mbox{练习~1.5.I}）。\footnote{译注：原书此处作「$\mathit{Mod}_{\OO_Y}$」，其中 $Y$ 未在本练习中出现，疑当为 $\mathfrak{U}$（\mbox{练习~2.7.D}／\mbox{练习~2.7.E} 中以 $Y$ 泛指开集），兹照录原书。——译者注}

\textbf{29.2.C.（练习）}\quad 设 $X$ 是一个\trRingedSpace，而 $\mathscr{G}$ 是一个 $\OO_X$-\trModule。

(a) 证明：
\[
\mathscr{E}xt^i(\OO_X,\mathscr{G})=
\begin{cases}
\mathscr{G}\text{，}& \text{若 } i=0\text{；}\\
0\text{，}& \text{其余情形。}
\end{cases}
\]

(b) 请描述一个\trCanonical\trIsomorphism{} $\mathrm{Ext}^i_X(\OO_X,\mathscr{G})\xleftarrow{\ \sim\ }\mathrm{H}^i(X,\mathscr{G})$。

\textbf{29.2.D.（练习）}\quad 用\mbox{练习~29.2.B} 与\mbox{练习~29.2.C(a)} 证明：若 $\mathscr{E}$ 是 $X$ 上的一个\trLocallyFreeSheaf，则当 $i>0$ 时 $\mathscr{E}xt^i(\mathscr{E},\mathscr{G})=0$。

在环上的\trModule{}的范畴中，我们喜欢\trProjectiveObj{}对象甚于内射对象，因为自由模易于打理。若\trScheme{}上的\trLocallyFreeSheaf{}总是\trProjectiveObj{}的，那就再好不过了，但遗憾的是这并不成立。因为若 $\OO_X$ 是投射的，那么当 $i>0$ 时 $\mathrm{Ext}^i(\OO_X,\cdot)$（作为函子）将为零；但由\mbox{练习~29.2.C(b)}，它就是函子 $\mathrm{H}^i(X,\cdot)$，而后者常常非零。尽管如此，我们仍然可以用一个局部自由分解来计算 $\mathscr{E}xt$，正如下面这道练习所示。

\textbf{29.2.E.（重要练习）}\quad 设 $X$ 是一个\trRingedSpace，而
\begin{equation*}
\tag{29.2.0.1}\label{eq:ext-lfr}
\cdots\longrightarrow\mathscr{E}_1\longrightarrow\mathscr{E}_0\longrightarrow\mathscr{F}\longrightarrow 0
\end{equation*}
是 $\OO_X$-\trModule{} $\mathscr{F}$ 的一个由\trLocallyFreeSheaf{}构成的分解。（当然，我们最感兴趣的场合是 $X$ 为\trScheme{}，而 $\mathscr{F}$ \trQC{}、甚至\trCoh{}的情形。）令 $\mathscr{E}_{\bullet}$ 表示 \eqref{eq:ext-lfr} 去掉 $\mathscr{F}$ 后所得的截断。请描述一个\trIsomorphism
\[
\mathscr{E}xt^i(\mathscr{F},\mathscr{G})\xleftarrow{\ \sim\ }\mathrm{H}^i\bigl(\mathscr{H}om(\mathscr{E}_{\bullet},\mathscr{G})\bigr).
\]
换言之，$\mathscr{E}xt^{\bullet}(\mathscr{F},\mathscr{G})$ 可以这样计算：取 $\mathscr{F}$ 的一个局部自由分解，加以截断，施加 $\mathscr{H}om(\cdot,\mathscr{G})$，再取同调。\footnote{译注：原书此处作「$\mathscr{E}xt^{\bullet}(\mathscr{F},\mathscr{G}))$」，右括号多印一枚，当为「$\mathscr{E}xt^{\bullet}(\mathscr{F},\mathscr{G})$」之误，径改。——译者注}提示：选取一个内射分解
\[
0\longrightarrow\mathscr{G}\longrightarrow\mathscr{I}_0\longrightarrow\mathscr{I}_1\longrightarrow\cdots
\]
并考虑以
\[
\begin{tikzcd}[row sep=2.2em, column sep=2.4em,font=\small]
\vdots & \vdots & & \\
\mathscr{H}om(\mathscr{E}_0,\mathscr{I}_1) \ar[r] \ar[u] & \mathscr{H}om(\mathscr{E}_1,\mathscr{I}_1) \ar[r] \ar[u] & \cdots \\
\mathscr{H}om(\mathscr{E}_0,\mathscr{I}_0) \ar[r] \ar[u] & \mathscr{H}om(\mathscr{E}_1,\mathscr{I}_0) \ar[r] \ar[u] & \cdots.
\end{tikzcd}
\]
为 $\mathrm{E}_0$ \trPage{}的\trSpectralSequence。

这个结果很重要：要计算 $\mathscr{E}xt$，我们可以用有限秩的局部自由分解来计算。你可以逐仿射地工作（借助\mbox{练习~29.2.B}），而在每一块仿射上，你可以对左变元使用自由分解。作为\mbox{练习~29.2.E} 的另一条推论：

\textbf{29.2.F.（练习）}

(a) 设 $\mathscr{F}$ 是\trLocally\trNoetherian{}\trScheme{} $X$ 上的一个\trCoherentSheaf，$\mathscr{G}$ 是 $X$ 上的一个\trQCSheaf。证明：$\mathscr{E}xt^i_X(\mathscr{F},\mathscr{G})$ 是 $X$ 上的一个\trQCSheaf。提示：由\mbox{练习~29.2.B}，由于 $\mathscr{E}xt$ 可以在\trOpenCover{}上计算，请证明只需对 $X$ \trAffine{}的情形证明该结果。但在这种情形下，$\mathscr{F}$ 有一个自由分解。（你将会发现 $\mathscr{E}xt^i_X(\mathscr{F},\mathscr{G})(\Spec A)=\mathrm{Ext}^i_A\bigl(\mathscr{F}(\Spec A),\mathscr{G}(\Spec A)\bigr)$。）

(b) 进而，若 $\mathscr{G}$ 也\trCoh，证明 $\mathscr{E}xt^i_X(\mathscr{F},\mathscr{G})$ 也是\trCoh{}的。

\textbf{29.2.1.}\ \emph{\trhRemark。}\quad 陈述「$\mathrm{Ext}^i(\mathscr{F},\mathscr{G})$ 与 $\mathscr{E}xt^i(\mathscr{F},\mathscr{G})$ 在第一个变元上的表现如同导出函子」在好几层意义上都成立。例如，我们有一条相应的\trLongExactSequence，见下一道练习。

\textbf{29.2.G.（练习）}\quad 设 $0\to\mathscr{F}'\to\mathscr{F}\to\mathscr{F}''\to 0$ 是\trRingedSpace{} $X$ 上诸 $\OO_X$-\trModule{}的一个\trExactSequence。对任意 $\OO_X$-\trModule{} $\mathscr{G}$，请描述一条\trLongExactSequence
\[
\begin{gathered}
0\longrightarrow\Hom(\mathscr{F}'',\mathscr{G})\longrightarrow\Hom(\mathscr{F},\mathscr{G})\longrightarrow\Hom(\mathscr{F}',\mathscr{G})\\
\longrightarrow\mathrm{Ext}^1(\mathscr{F}'',\mathscr{G})\longrightarrow\mathrm{Ext}^1(\mathscr{F},\mathscr{G})\longrightarrow\mathrm{Ext}^1(\mathscr{F}',\mathscr{G})\longrightarrow\cdots.
\end{gathered}
\]
（于是，$\mathrm{Ext}^i(\cdot,\mathscr{G})$ 构成一个反变\trDeltaFunctor——如 \S23.2.7 中所定义——尽管它并不是作为导出函子来定义的。）请对 $\mathscr{E}xt$ 证明类似的结果。提示：取一个内射分解 $0\to\mathscr{G}\to\mathscr{I}^0\to\cdots$。利用「若 $\mathscr{I}$ 内射，则 $\Hom(\cdot,\mathscr{I})$ \trExact」（内射性的定义，\mbox{练习~23.2.C(a)}）这一事实。由此得到\trComplexN{}的一个\trShortExactSequence
\[
0\longrightarrow\Hom(\mathscr{F}'',\mathscr{I}^{\bullet})\longrightarrow\Hom(\mathscr{F},\mathscr{I}^{\bullet})\longrightarrow\Hom(\mathscr{F}',\mathscr{I}^{\bullet})\longrightarrow 0
\]
并取上同调中的\trLongExactSequence。

\textbf{29.2.2. 对 $\mathrm{Ext}$ 的函子性 \trSerreDuality{}在\trProjectiveSpace{}上成立。}

我们现在证明：对 $\mathrm{Ext}$ 的函子性 \trSerreDuality{}在\trProjectiveSpace{}上成立。我们将使用泛\trDeltaFunctor{}的机制（\S23.2.6），所以你既不妨快速浏览一下那一节，也可以干脆忽略这段讨论。

\textbf{29.2.H.（练习）}\quad 证明：$\bigl(\mathrm{Ext}^i_{\mathbb{P}^n_k}(\cdot,\omega_{\mathbb{P}^n_k})\bigr)$ 是一个（反变的）泛\trDeltaFunctor。提示：$\mathrm{Ext}$ 在其第一个变元上不是导出函子，所以你不能用\mbox{推论~23.2.11} 的「\trProjectiveObj」版本。作为替代，请使用\mbox{定理~23.2.9}，以及「对每个 $\mathscr{F}$ 都存在满态射 $\OO(m)^{\oplus N}\to\mathscr{F}$，其中 $m<0$」这一事实。

\textbf{29.2.I.（练习）}\quad 证明：$\bigl(\mathrm{H}^{n-i}(\mathbb{P}^n_k,\cdot)^{\vee}\bigr)$ 是一个泛\trDeltaFunctor。（诸 $\delta$-映射是什么？）提示：试用与上一道练习相同的想法。

\mbox{命题~29.1.10} 给出函子之间的一个\trIsomorphism{} $\mathrm{Ext}^0_{\mathbb{P}^n_k}(\cdot,\omega_{\mathbb{P}^n_k})\xleftarrow{\ \sim\ }\mathrm{H}^n(\mathbb{P}^n_k,\cdot)^{\vee}$，于是由泛\trDeltaFunctor{}的\mbox{定义~23.2.8}，我们便有诸\trDeltaFunctor{}之间的一个\trIsomorphism{} $\bigl(\mathrm{Ext}^i_{\mathbb{P}^n_k}(\cdot,\omega_{\mathbb{P}^n_k})\bigr)\xleftarrow{\ \sim\ }\bigl(\mathrm{H}^{n-i}(\mathbb{P}^n_k,\cdot)^{\vee}\bigr)$，从而证明了对 $\mathbb{P}^n_k$ 而言对 $\mathrm{Ext}$ 的函子性 \trSerreDuality。

\textbf{29.2.3. 两道有用的练习，以及\trProjectiveSpace{}上对\trVectorBundle{}的函子性 \trSerreDuality。}

\textbf{29.2.J.（练习）}\quad 设 $X$ 是一个\trRingedSpace，$\mathscr{F}$ 与 $\mathscr{G}$ 是 $\OO_X$-\trModule，而 $\mathscr{E}$ 是 $X$ 上的一个\trLocallyFreeSheaf。请描述诸\trIsomorphism
\[
\mathscr{E}xt^i(\mathscr{F}\otimes\mathscr{E}^{\vee},\mathscr{G})\xleftarrow{\ \sim\ }\mathscr{E}xt^i(\mathscr{F},\mathscr{G}\otimes\mathscr{E})\xleftarrow{\ \sim\ }\mathscr{E}xt^i(\mathscr{F},\mathscr{G})\otimes\mathscr{E}
\]
以及
\[
\mathrm{Ext}^i(\mathscr{F}\otimes\mathscr{E}^{\vee},\mathscr{G})\xleftarrow{\ \sim\ }\mathrm{Ext}^i(\mathscr{F},\mathscr{G}\otimes\mathscr{E}).
\]
提示：证明若 $\mathscr{I}$ 内射，则 $\mathscr{I}\otimes\mathscr{E}$ 也内射。

\textbf{29.2.K.（练习）}\quad 若 $\mathscr{F}$ 是\trScheme{} $X$ 上的一个\trLocallyFreeSheaf，而 $\mathscr{G}$ 是 $X$ 上任意一个\trCoherentSheaf，请描述一个\trIsomorphism{} $\mathrm{Ext}^i(\mathscr{F},\mathscr{G})\xleftarrow{\ \sim\ }\mathrm{H}^i(X,\mathscr{F}^{\vee}\otimes\mathscr{G})$（对 $\mathscr{F}$ 函子）。取 $\mathscr{G}$ 为 $\omega_X$，证明：对 $\mathrm{Ext}$ 的函子性 \trSerreDuality{}蕴含对\trVectorBundle{}的函子性 \trSerreDuality。（你也不妨琢磨一下迹版本。）于是，由 \S29.2.2，对\trVectorBundle{}的函子性 \trSerreDuality{}在\trProjectiveSpace{}上成立。提示：\mbox{练习~29.2.J} 与\mbox{练习~29.2.C(b)}。

\textbf{29.2.4. $\mathrm{Ext}$ 的局部到整体\trSpectralSequence。}

「层」$\mathscr{E}xt$ 与「整体」$\mathrm{Ext}$ 之间由一个\trSpectralSequence{}相联系。一旦我们证明了 $\mathscr{H}om(\mathscr{F},\mathscr{I})$ 对函子 $\Gamma$ 是\trAcyclic{}的，这就是\trGrothendieck{}函子复合谱序列（\mbox{定理~23.3.5}）的直接应用。

\textbf{29.2.L.（练习）}\quad 设 $\mathscr{I}$ 是一个内射 $\OO_X$-\trModule。证明：$\mathscr{H}om(\mathscr{F},\mathscr{I})$ 是\trFlasque{}的（\mbox{定义~23.4.6}），从而由\mbox{练习~23.4.I} 是 $\Gamma$-\trAcyclic{}的。提示：设 $j\colon\mathfrak{U}\hookrightarrow V$ 是\trOpenSubset{}之间的一个包含。我们希望证明 $\mathscr{H}om(\mathscr{F},\mathscr{I})(V)\to\mathscr{H}om(\mathscr{F},\mathscr{I})(\mathfrak{U})$ 是\trSurjectiveMap。注意 $\mathscr{I}|_V$ 在 $V$ 上内射（\mbox{练习~23.5.A}）。把\trExact{}函子 $\Hom_V(\cdot,\mathscr{I}|_V)$ 施加到 $V$ 上诸\trSheaf{}之间的包含态射 $j_!(\mathscr{F}|_{\mathfrak{U}})\hookrightarrow\mathscr{I}|_V$ 上（\S23.4.7）。

\textbf{29.2.M.（练习：$\mathrm{Ext}$ 的局部到整体\trSpectralSequence）}\quad 设 $X$ 是一个\trRingedSpace，而 $\mathscr{F}$ 与 $\mathscr{G}$ 是 $\OO_X$-\trModule。请描述一个以 $\mathrm{H}^i\bigl(X,\mathscr{E}xt^j(\mathscr{F},\mathscr{G})\bigr)$ 为 $\mathrm{E}_2$ 项、\trAbut{}于 $\mathrm{Ext}^{i+j}(\mathscr{F},\mathscr{G})$ 的\trSpectralSequence。（提示：使用\trGrothendieck{}函子复合谱序列，\mbox{定理~23.3.5}。回忆 $\Hom(\mathscr{F},\cdot)=\Gamma\bigl(\mathscr{H}om(\mathscr{F},\cdot)\bigr)$，\mbox{练习~2.3.C}。）
